\section{The Greater Mysteries}

\begin{quotex}
Omnipotence does not mean power to do absurdities. The compulsion of another's will is such an absurdity, and therefore no real omnipotence could force such a compulsion. Omnipotence is spiritual, and spirit acts not by brute compulsion but by knowledge and inspiration. \flright{\textsc{R. G. Collingwood}}

\end{quotex}
Over the past 150 years or so, we have been gifted with some extraordinary revelations which conceivably will buttress the foundation of a spiritual renewal. Apart from some small pockets of interest, the world continues on its downward path. If you are happy with your current spiritual reading, perhaps these revelations don't matter much. If you are entrenched in materialism and scientism, these ideas may not be able to penetrate your minds. But for a few of us, these ideas form the foundation for the worldview of Tradition.

To commence, I shall mention these three in particular:

\begin{itemize}
\item \textbf{Vladimir Solovyov} revealed the inner connection between Egyptian theosophy and Christian doctrine. 
\item \textbf{Rene Guenon} revealed the existence of the Primordial Tradition. 
\item \textbf{Georges Dumezil} revealed the social structure common to traditional societies. 
\end{itemize}
\paragraph{Hermetic Roots of Tradition}
In an intriguing footnote to his most metaphysical book, Lectures on Divine Humanity, Vladimir Solovyov noted the close connection between Hermetic philosophy and Christian doctrine. Specifically, the ideas of the Logos and the Trinity was known to the Neoplatonic philosophers. His conclusion was that Christian doctrine did not create a new metaphysics, but its originality is in its revelation of positive facts. Viz., the Logos was incarnated in Jesus, and the Trinity included three divine persons, beyond its abstract metaphysical principles.

There are those who regard the reawakening of Hermetic ideas in the Renaissance as some sort of aberration and the sign of a decline. Rather, it is a recovery of what was there in the beginning. The decline is because of the separation of Tradition from the exoteric organisation. The result is that the esoteric teaching veered off into phantasies and the exoteric lost its soul, and has been carried about with every wind of doctrine by the wickedness of men.

In our own time, we have the works of \textbf{Valentin Tomberg} and \textbf{Boris Mouravieff} who, in different ways, have reunited the two streams.

\paragraph{The Primordial Tradition}
Despite all of his works being readily available, Rene Guenon is still not properly understood. Some say he created a novel ``School" of ``Traditionalism" with certain beliefs and dogmas. To the contrary, he revealed the existence of the Primordial Tradition, which is not a matter of speculation.

The source of the Primordial Tradition is superhuman, and its metaphysical teaching is eternal. Hence, it is not contingent on any human activities and moreover it has always existed and will always exist. A fortiori, it is not the teaching of a particular ``School". Nor can one be a vague ``traditionalist"; it is necessary to be part of a specific tradition.

Nevertheless, the Tradition does take form in various cultures and historical epochs. It is then expressed in doctrines, rites, and symbols appropriate to its exoteric manifestation. These manifestations of the Primordial Tradition also show the way to the mysteries:

\begin{itemize}
\item The Lesser Mysteries show the way back to the primordial state before the Fall of man 
\item The Greater Mysteries show the path to gnosis through self-knowledge 
\end{itemize}
The primordial state is the natural state of humanity, even if its memory has been mostly lost. Since we have already, in multiple posts, discussed the Primordial Tradition unveiled by Guenon, there is not need to repeat it all here.

\paragraph{Trifunctional Social Organisation}
In his cross-cultural studies of myths and social institutions, Georges Dumezil discovered the tri-functional organisation of Indo-European cultures. These are:

\begin{itemize}
\item The function of the sacred 
\item The warrior function 
\item The producers 
\end{itemize}
There correspond to three human types:

\begin{itemize}
\item Those who are strong by intelligence 
\item Those who are strong in courage and valour in battle 
\item Those who are wealthy 
\end{itemize}
In the European Middle Ages, society was similarly organised:

\begin{itemize}
\item Oratores: those who pray (the clergy) 
\item Bellatores: those who fight (the nobility) 
\item Laboratores: those who work 
\end{itemize}
Moreover, these three functions were hierarchically organised, so that the producers were under the nobility or the political power, who themselves were subject to the spiritual authority. There was no freedom of religion, since the whole society was based on spiritual unity. To be excommunicated was tantamount to exile from the group. Modern theories, on the other hand, try to achieve unity based on other contingent factors such as race, economic class, or political party. Such societies are fragile, whereas societies based on a common spiritual foundation have tended to persist.

\paragraph{Acres of Diamonds}
In \emph{The Vedanta and Western Tradition}\footnote{\url{https://gornahoor.net/?p=556}}, \textbf{Ananda Coomaraswamy} listed eight Western thinkers who, in their totality, come closest to the teachings of the Vedanta. These are, at a minimum:

\begin{enumerate}
\item Plato 
\item Philo 
\item Hermes Trismegistus 
\item Plotinus 
\item Gospel of John 
\item Dionysius the Areopagite 
\item Meister Eckhart 
\item Dante 
\end{enumerate}
I would add four more:

\begin{enumerate}
\item Aristotle 
\item Augustine of Hippo 
\item Boethius 
\item Thomas Aquinas 
\end{enumerate}
With that grounding, there can be found:

\begin{itemize}
\item Knowledge of the Absolute 
\item Cosmology 
\item The distinction between the worlds of Being and Becoming 
\item Hylomorphism 
\item The psychological constitution of human beings 
\item An ethical system based on virtues and excellence 
\item Political theory 
\item The three stages of spiritual life, and so on 
\end{itemize}
Yet, of those claiming to be ``traditional", who takes this list seriously? Instead, we hear pointless arguments about what is tradition. For example, is European tradition pagan or Christian? If you decide on the former, then this list still applies to you; or else you are just role playing.

Other Westerners decide to look elsewhere, becoming Druids, Vedantic priests, Buddhists, etc. However, if they cannot understand their own tradition, can you really trust that they can understand another?

\textbf{Russell Conwell} wrote the famous essay \textit{Acres of Diamonds}\footnote{\url{https://americanrhetoric.com/speeches/rconwellacresofdiamonds.htm}} about a man who went off in search of diamonds, when there were diamonds scattered on his own property that he abandoned. That is the situation we are in today, always searching the next glittery thing elsewhere, instead of where it is actually found.

\paragraph{The Absolute and God}
The strictly philosophical point of view can only rise to the level of the Absolute. Starting with the realisation that all things and events are interrelated, one rises up to an understanding that encompasses increasing number of things in space and time. Ultimately, one is led to the Absolute. However, different thinkers have understood it in different way; for example:

\begin{itemize}
\item Plato: the form of the Good 
\item Aristotle: the unmoved mover and first cause 
\item Plotinus: the One 
\item Spinoza: Nature 
\item Schopenhauer: Blind Will 
\item Francis Bradley: the True 
\end{itemize}
Some philosophers believe that mankind is evolving toward the Absolute, while in most other systems, the notion of change is an illusion.

Hence, philosophy is part of a larger conversation. Sharing in that has its positive points. The intellectual contemplation of the Absolute can be accompanied by joy and peace. However, this is at the level of the lesser mysteries.

If the human being is a person, the question must needs arise about the Absolute: How can an abstract Absolute give rise to persons with their own interiority? The Absolute itself, then, must also have a centre of interiority. Once this is grasped, then the idea of revelation from, and relationship to, the Absolute makes sense. God reveals himself as the Absolute, yet is incomprehensible. An object can be known, but a subject, or person, always contains something hidden.

At the stage of the contemplation of God, unlike philosophical contemplation of the Absolute, there is no subject-object difference. There is just one Self. That is the goal of the greater mysteries.

\paragraph{The European Tradition}
A final point for those who haven't yet grasped the idea of Tradition. Since it is superhuman and transcendental, it does not depend on individual men and women. If the spiritual leaders fail, if their organisation fails, the Tradition itself perdures. This is timely, given the confused state of the organisation whose duty is to preserve the Western tradition. Keep this in mind:

\begin{quotex}
The Church is not a man or even a mere collective of men, but the Mystical Body of Christ whose subsistence cannot be destroyed by any human error. Lost in the current mania of hyper-papalism is the infallibility of the Church as a corporate whole, extending even to the faithful as a body, which obeys what has always been taught by the Church as a whole and rejects what is foreign to that teaching. \flright{\emph{The Remnant Newspaper}, March 31, 2019.}

\end{quotex}
In other words, follow tradition, not the organisation.



\flrightit{Posted on 2019-06-26 by Cologero }

\begin{center}* * *\end{center}

\begin{footnotesize}\begin{sffamily}



\texttt{Sam Viehweg on 2019-06-27 at 18:31 said: }

Can I ask, regarding the following: ``Nor can one be a vague ``traditionalist"; it is necessary to be part of a specific tradition."

Where does this leave Evola? And where does this leave one who has no avenue externally to be part of a specific tradition?


\hfill

\texttt{Cologero on 2019-06-27 at 18:56 said: }

We are in a mess

\url{https://www.youtube.com/watch?v=AjplZXgodhs}


\hfill

\texttt{Tom on 2019-07-10 at 11:59 said: }

Why is it automatically assumed that being is superior to becoming ?


\hfill

\texttt{Cologero on 2019-07-13 at 23:09 said: }

``superior" simply means higher up. Our world of experience is the world of becoming. The modern world automatically assumes that this world is causally closed, thereby failing to recognize transcendence. Unfortunately, the sensory world of experience is unintelligible. This is too obvious to waste much time on it. Viz., people may agree on the facts of experience, but there is no agreement on the meaning of those facts. Surely you have noticed that. In Platonic terms, we can only have opinions about the world of becoming, but not knowledge.

Real knowledge is knowledge of essences, which are part of the world of being. I've spent a lot of virtual ink illustrating that in many ways, from many angles. Pick a particular sentence or paragraph, if you have a question.


\end{sffamily}\end{footnotesize}
